%% start of file `template.tex'.
%% Copyright 2006-2013 Xavier Danaux (xdanaux@gmail.com).
%
% This work may be distributed and/or modified under the
% conditions of the LaTeX Project Public License version 1.3c,
% available at http://www.latex-project.org/lppl/.
\documentclass[11pt,a4paper,roman,colorlinks,linkcolor = blue]{moderncv}        % possible options include font size ('10pt', '11pt' and '12pt'), paper size ('a4paper', 'letterpaper', 'a5paper', 'legalpaper', 'executivepaper' and 'landscape') and font family ('sans' and 'roman')


% modern themes
\moderncvstyle{banking}                            % style options are 'casual' (default), 'classic', 'oldstyle' and 'banking'
\moderncvcolor{black}                                % color options 'blue' (default), 'orange', 'green', 'red', 'purple', 'grey' and 'black'
%\renewcommand{\familydefault}{\sfdefault}         % to set the default font; use '\sfdefault' for the default sans serif font, '\rmdefault' for the default roman one, or any tex font name
\nopagenumbers{}                                  % uncomment to suppress automatic page numbering for CVs longer than one page

% character encoding
%\usepackage[utf8]{inputenc}

\usepackage{fontspec}
\usepackage{tabularx}
\usepackage{ragged2e}
\usepackage{url}

% if you are not using xelatex ou lualatex, replace by the encoding you are using
%\usepackage{CJKutf8}                              % if you need to use CJK to typeset your resume in Chinese, Japanese or Korean

% adjust the page margins
\usepackage[scale=0.8]{geometry}
\usepackage{multicol}
%\setlength{\hintscolumnwidth}{3cm}                % if you want to change the width of the column with the dates
%\setlength{\makecvtitlenamewidth}{10cm}           % for the 'classic' style, if you want to force the width allocated to your name and avoid line breaks. be careful though, the length is normally calculated to avoid any overlap with your personal info; use this at your own typographical risks...

\usepackage{import}
\usepackage{apacite}
\newcommand*{\doi}[1]{\href{http://dx.doi.org/#1}{#1}}
% personal data
\name{Thom Benjamin}{Volker}
% \title{Curriculum Vitae}                               % optional, remove / comment the line if not wanted

\address{Utrecht, The Netherlands}{}{}% optional, remove / comment the line if not wanted; the "postcode city" and and "country" arguments can be omitted or provided empty

% \phone[mobile]{909-839-3097}                   % optional, remove / comment the line if not wanted
% \phone[fixed]{01234 123456}                    % optional, remove / comment the line if not wanted
%\phone[fax]{+3~(456)~789~012}                      % optional, remove / comment the line if not wanted
% \email{xpan1@swarthmore.edu}                               % optional, remove / comment the line if not wanted
% \homepage{shawnpan.me}                         % optional, remove / comment the line if not wanted
% \extrainfo{}                 % optional, remove / comment the line if not wanted
%\photo[64pt][0.4pt]{picture}                       % optional, remove / comment the line if not wanted; '64pt' is the height the picture must be resized to, 0.4pt is the thickness of the frame around it (put it to 0pt for no frame) and 'picture' is the name of the picture file
%\quote{Some quote}                                 % optional, remove / comment the line if not wanted

% to show numerical labels in the bibliography (default is to show no labels); only useful if you make citations in your resume
%\makeatletter
%\renewcommand*{\bibliographyitemlabel}{\@biblabel{\arabic{enumiv}}}
%\makeatother
%\renewcommand*{\bibliographyitemlabel}{[\arabic{enumiv}]}% CONSIDER REPLACING THE ABOVE BY THIS

% bibliography with mutiple entries
%\usepackage{multibib}
%\newcites{book,misc}{{Books},{Others}}
  
\newcommand*{\customcventry}[7][.25em]{
  \begin{tabular}{@{}l} 
    {\bfseries #4}
  \end{tabular}
  \hfill% move it to the right
  \begin{tabular}{l@{}}
     {\bfseries #5}
  \end{tabular} \\
  \begin{tabular}{@{}l} 
    {\itshape #3}
  \end{tabular}
  \hfill% move it to the right
  \begin{tabular}{l@{}}
     {\itshape #2}
  \end{tabular}
  \ifx&#7&%
  \else{\\%
    \begin{minipage}{\maincolumnwidth}%
      \small#7%
    \end{minipage}}\fi%
  \par\addvspace{#1}}

\newcommand*{\customcvproject}[4][.25em]{
%   \vfill\noindent
  \begin{tabular}{@{}l} 
    {\bfseries #2}
  \end{tabular}
  \hfill% move it to the right
  \begin{tabular}{l@{}}
     {\itshape #3}
  \end{tabular}
  \ifx&#4&%
  \else{\\%
    \begin{minipage}{\maincolumnwidth}%
      \small#4%
    \end{minipage}}\fi%
  \par\addvspace{#1}}

\setlength{\tabcolsep}{12pt}

%----------------------------------------------------------------------------------
%            content
%----------------------------------------------------------------------------------
\begin{document}

\definecolor{darkblue}{HTML}{153884}
\hypersetup{urlcolor=darkblue}
%\begin{CJK*}{UTF8}{gbsn}                          % to typeset your resume in Chinese using CJK
%-----       resume       ---------------------------------------------------------
\makecvtitle
\vspace*{-13mm}

\begin{center}
\begin{tabular}{ c c c c }
 \faGlobe\enspace \href{https://thomvolker.github.io}{thomvolker.github.io} & \faEnvelope\enspace \href{mailto:t.b.volker@uu.nl}{t.b.volker@uu.nl} & \faGithub\enspace \href{https://www.github.com/thomvolker}{GitHub} & \faLinkedin\enspace \href{https://www.linkedin.com/in/thom-volker-a4620415a/}{LinkedIn}% & \faMobile\enspace 123-456-7891\\
\end{tabular}
\end{center}

%----------------------------------------------------------------------------------
%            EDUCATION
%----------------------------------------------------------------------------------

\section{Education and experience}

{\customcventry{}{Utrecht University, Statistics Netherlands}{PhD Candidate Synthetic Data}{2022 - Current}{}{}
{\begin{itemize}
  \item[$\circ$] Supervisors: S. van Buuren, E.-J. van Kesteren, P.-P. de Wolf
\end{itemize}
}
}

{\customcventry{Valedictorian, Cum laude (GPA 8.9/10)}{Utrecht University}{MSc Methodology and Statistics}{2019 - 2022}{}{}
{\begin{itemize}
  \item[$\circ$] Internship missing data with Gerko Vink and Stef van Buuren
  \item[$\circ$] Various student-assistant jobs: data visualization, Shiny app development, contributions to \href{https://asreview.nl/}{ASReview} and teaching.
\end{itemize}
}
}
\customcventry{Cum laude (GPA 8.7/10)}{Utrecht University}{MSc Sociology and Social Research}{2019 - 2022}{}{}

\customcventry{(GPA 7.7/10)}{Utrecht University}{BA Liberal Arts \& Sciences}{2016 - 2019}{}{}{}


%----------------------------------------------------------------------------------
%            WORKING EXPERIENCE
%----------------------------------------------------------------------------------



\section{Teaching}
\begin{itemize}
\item[$\circ$] \textbf{Post graduate-level courses}
{\begin{itemize}
  \item[--] \href{https://utrechtsummerschool.nl/courses/data-science/data-science-solving-missing-data-problems-in-r}{Solving Missing Data Problems in \texttt{R}} [2019 - 2023]: \textit{Developing materials, supervising practicals.}
  \item[--] \href{http://giachanou.com/machine-learning/}{Machine learning with python} [2024]: \textit{Supervising practicals.}
  \item[--] \href{https://utrechtsummerschool.nl/courses/data-science/data-science-statistical-programming-with-r}{Statistical Programming with R} [2021]: \textit{Supervising practicals.}
  \item[--] \href{https://utrechtsummerschool.nl/courses/data-science/survey-research-statistical-analysis-and-estimation}{Survey Research: Statistical Analysis and Estimation} [2021]: \textit{Supervising practicals.}
\end{itemize}
  }
  \item[$\circ$] \textbf{Master level courses}
  {\begin{itemize}
    \item[--] Data wrangling [2021 - Current]: \textit{Supervising practicals and grading exams.} 
    \item[--] Battling the curse of dimensionality [2022 - 2025]: \textit{Developing course materials, supervising practical sessions and grading exams.}
    \item[--] Network Analysis [2020 - 2021]: \textit{Developing course materials, supervising practicals and grading exams.}
    \item[--] Methodological and Statistical Aspects of Social Science Research [2019 - 2020]: \textit{Supervising practicals.}
   \end{itemize}
  }
  \item[$\circ$] \textbf{Bachelor level courses}
  {\begin{itemize}
    \item[--] Theory Construction and Statistical Modeling [2019- 2020]: \textit{Supervising practicals}
    \item[--] Various undergraduate courses on standard statistical methods [2018-2021].
   \end{itemize}
  }
\end{itemize}

\section{Supervision}
\begin{itemize} 
  \item[$\circ$] Ludo Mels (Master's thesis, co-supervised with \href{https://www.saravanerp.com/}{Dr. Sara van Erp}, 2026-2027).
  \item[$\circ$] Heleen Br{\"u}ggen (Master's thesis, co-supervised with \href{https://hanneoberman.github.io/}{Hanne Oberman} \& \href{https://www.gerkovink.com}{Gerko Vink}, 2023-2024).
  \item[$\circ$] Mirthe Hendriks (Master's thesis, co-supervised with \href{https://www.gerkovink.com}{Gerko Vink}, 2021).
  \item[$\circ$] Stijn van den Broek (Master's thesis, co-supervised with \href{https://www.gerkovink.com}{Gerko Vink}, 2021). 
  \item[$\circ$] Romain Rey (Bachelor's thesis, co-supervised with \href{https://www.uu.nl/medewerkers/iklugkist}{Irene Klugkist}, 2020). 
\end{itemize}


%----------------------------------------------------------------------------------
%            VOLUNTARY WORK AND OTHER THINGS
%----------------------------------------------------------------------------------

\section{Miscellanea}

\textbf{IOPS PhD Representative [2022 - 2024]:} Represent PhD candidates within the Interuniversity Graduate School of Psychometrics and Sociometrics.
\textbf{Organization Linear Algebra Reading Group [2025-2026]}
\textbf{Organization Human Data Science Reading Group [2022-2024]}
\textbf{Social committee [2022 - Current]:} Organising frequent social activities for the department. 
\textbf{Debuut [2019 - 2020]:} Participate in buddy programme at Utrecht University to match potential first generation university students ($\sim$ 11 year olds) with actual students to get a flavour of what studying at a university entails.

%----------------------------------------------------------------------------------
%            SOFTWARE PROFICIENCY
%----------------------------------------------------------------------------------

\section{Received funding}

{\customcventry{}{Project lead}{SynthEval}{2023-2024}{}
{\begin{itemize}
  \item[$\circ$] An R package for evaluating and improving the quality of synthetic data sets. For €15.000.
\end{itemize}
  }
}



% Publications from a BibTeX file without multibib
%  for numerical labels: \renewcommand{\bibliographyitemlabel}{\@biblabel{\arabic{enumiv}}}% CONSIDER MERGING WITH PREAMBLE PART

\section{Peer-reviewed publications}

\begin{itemize}
\item[$\circ$] Klugkist, I. \& \textbf{Volker, T. B.} (2023). Bayesian Evidence Synthesis for Informative Hypotheses: An introduction. \textit{Psychological Methods, Advance online publication.} \href{https://psycnet.apa.org/fulltext/2024-05522-001.html}{Paper}.
\item[$\circ$] Van Leeuwen, A., Bagheri, A., \textbf{Volker, T. B.} \& Van Brakel, Charlotte (2023). Verkenning van de inzet van topic modelling bij het analyseren van schrijfopdrachten in de context van de flipped classroom. [Exploring the use of topic modelling for the analysis of written assignments in a flipped classroom context.] \textit{Tijdschrift voor Hoger Onderwijs, 41} (2), 84-98. \href{https://tvho.nl/article/download/15690/17434}{Paper}.
\item[$\circ$] \textbf{Volker, T. B.} \& Vink, G. (2021). Anony\textit{mice}d Shareable Data: Using \texttt{mice} to Create and Analyze Multiply Imputed Synthetic Datasets. \textit{psych, 3}, 703-716. \href{https://www.mdpi.com/2624-8611/3/4/45}{Paper}.
\item[$\circ$] \textbf{Volker, T. B.}, Buskens, V. \& Raub, W. (Forthcoming). How Networks Facilitate Trust: A Synthesis on Control Effects. In \textit{Handbook of Trust \& Social Networks}.
\end{itemize}

\section{Software development}

\begin{itemize}
\item[$\circ$] \texttt{densityratio} [aut] - Developed \texttt{R}-package for directly estimating the ratio beteen two sample densities.
\item[$\circ$] \texttt{mice} [ctb] - Various additions to the multiple imputation software \texttt{mice}.
\item[$\circ$] \texttt{jaspMissingData} [ctb] - A JASP module for multiple imputation with the \texttt{R}-package \texttt{mice}.
\item[$\circ$] \texttt{ggmice} [ctb] - Adjusted functionality to allow for visualization of synthetic data.
\end{itemize}

\section{Preprints, vignettes, blogs and other writings}

\begin{itemize}
\item[$\circ$] \textbf{Volker, T.B.}, De Wolf, P.-P. \& Van Kesteren, E. J. (Under review). A density ratio framework for evaluating the utility of synthetic data. \href{https://arxiv.org/pdf/2408.13167}{Paper}.
\item[$\circ$] \textbf{Volker, T.B.}. Generating and evaluating synthetic data in \texttt{R}. \href{https://lmu-osc.github.io/synthetic-data-tutorial/}{Commissioned tutorial for LMU Open Science Center}.
\item[$\circ$] Poses, C. \& \textbf{Volker, T.B.} (2025). \texttt{densityratio}: An \texttt{R}-package for distribution comparison through density ratio estimation. \href{https://thomvolker.github.io/densityratio/articles/densityratio.html}{Software vignette}.
\item[$\circ$] Poses, C. \& \textbf{Volker, T.B.} (2025). Covariate shift adjustment with the \texttt{densityratio} \texttt{R}-package. \href{https://thomvolker.github.io/densityratio/articles/covariate-shift.html}{Software vignette}.
\item[$\circ$] \textbf{Volker, T.B.} (2025). High dimensional two-sample testing with the \texttt{densityratio} \texttt{R}-package. \href{https://thomvolker.github.io/densityratio/articles/high-dim-testing.html}{Software vignette}.
\item[$\circ$] \textbf{Volker, T.B.} (2025). Illustrative GANs for simply synthetic data with \texttt{keras} in \texttt{R}. \href{https://thomvolker.github.io/blog/2507_gans_in_r/}{Blog post}.
\item[$\circ$] \textbf{Volker, T.B.} (2025). Different ways of calculating OLS regression coefficients (in \texttt{R}). \href{https://thomvolker.github.io/blog/2506_regression/}{Blog post}.
\item[$\circ$] \textbf{Volker, T.B.} (2025). Statistical power of the Bayesian one-sample t-test. \href{https://thomvolker.github.io/bf-power-ttest/bayes-ttest-power.html}{Blog post}.
\item[$\circ$] Research Data Management Support (2023). Data Privacy Handbook. \href{https://utrechtuniversity.github.io/dataprivacyhandbook/}{Online book}.
\item[$\circ$]\textbf{Volker, T.B.} \& Klugkist, I. (2023). Combining support for hypotheses over heterogeneous studies with Bayesian Evidence Synthesis: A simulation study. \href{https://arxiv.org/pdf/2312.15032.pdf}{arXiv}.
\item[$\circ$] \textbf{Volker, T.B.} \& Vink, G. (2022). \texttt{futuremice}: The future starts today. \url{https://www.gerkovink.com/futuremice/Vignette_futuremice.html}{Software vignette}.
\end{itemize}


\section{Invited workshops, talks and conference presentations}
\begin{itemize}
\item[$\circ$] Synthetic data: How fake data can be useful. Guest lecture for elective course \textit{Mimicking \& Masking, Utrecht University}, 8 November 2026. \href{https://thomvolker.github.io/masking-lecture/}{Slides}.
\item[$\circ$] \texttt{densityratio}: An \texttt{R}-package for distribution comparison. Talk for ETH Z{\"u}rich, Switzerland, 30 April 2026, . \href{https://thomvolker.github.io/densityratio-zuerich/}{Slides}.
\item[$\circ$] Synthetic data for Enhancing open and reproducible social science. Invited talk for UTH Z{\"u}rich, Switzerland, 23 April 2026. \href{https://thomvolker.github.io/synzuerich/}{Slides}.
\item[$\circ$] A practical introduction to multiple imputation of missing data with the \texttt{R}-package \texttt{mice}. Invited workshop for Workshops for Ukraine, 22 January 2026. \href{https://thomvolker.github.io/mice4ukraine}{Slides}.
\item[$\circ$] Statistical disclosure control for medical output. Invited talk for VAC4EU Consortium, 11 November 2025. \href{https://thomvolker.github.io/sdc4vac4eu/}{Slides}.
\item[$\circ$] Missing data in medical research: From practical problems to flexible solutions [with Hanne Oberman], Invited half day course at St. Antonius Hospital, 9 November 2025. \href{https://hanneoberman.github.io/antonius/00_index.html}{Course materials}.
\item[$\circ$] Prediction intervals with missing data [with Florian van Leeuwen and Stef van Buuren], Conference presentation at ODISSEI, 4 November 2025. \href{https://thomvolker.github.io/predint-odissei/predint-odissei.html}{Slides}.
\item[$\circ$] Synthetic data: A solution for data access problems [with Peter Gerbrands, in Dutch], Invited talk at the Dutch Chamber of Commerce, 20 November 2025. \href{https://thomvolker.github.io/talks/presentations/kvk_synthetic_data.pdf}{Slides}.
\item[$\circ$] How networks facilitate trust: A synthesis on control effects [with Werner Raub and Vincent Buskens]. Conference presentation at Trust and Social Networks conference, 23 July 2025. \href{https://osf.io/pw7z5/files/2xdy9}{Slides}.
\item[$\circ$] Measuring utility of synthetic data [in Dutch]. Invited talk for Dutch Knowledge Network Synthetic Data, 19 December 2023. \href{https://thomvolker.github.io/utiliteit/}{Slides}.
\item[$\circ$] Assessing the utility of synthetic data: A density ratio perspective [with Peter-Paul de Wolf \& Erik-Jan van Kesteren], UNECE Expert Meeting on Statistical Data Confidentiality, 26-28 September 2023, Wiesbaden, Germany. \href{https://github.com/thomvolker/unece-density-ratio/tree/main/presentation}{Slides}.
\item[$\circ$] Fake it ‘till you make it: Generating synthetic data with high utility in \texttt{R} [with Erik-Jan van Kesteren], T{\"u}bingen Open Science Winter School, 13 February 2023, T{\"u}bingen, Germany. \href{https://thomvolker.github.io/OSWS_Synthetic/}{Course materials}.
\item[$\circ$] Creating synthetic data with \texttt{mice} in \texttt{R} [with Gerko Vink], UMCU, 4 November2022, Utrecht, the Netherlands. \href{https://www.gerkovink.com/syn/}{Slides}.
\end{itemize}


%-----       letter       ---------------------------------------------------------

\end{document}


%% end of file `template.tex'.
